%!TeX root=../tese.tex
%("dica" para o editor de texto: este arquivo é parte de um documento maior)
% para saber mais: https://tex.stackexchange.com/q/78101

%% ------------------------------------------------------------------------- %%

% "\chapter" cria um capítulo com número e o coloca no sumário; "\chapter*"
% cria um capítulo sem número e não o coloca no sumário. Alguns autores
% preferem que a introdução não seja numerada, mas ainda assim apareça
% no sumário. Para isso, você pode usar "\chapter**{Introdução}". Esse
% comando não é padrão do LaTeX; ele foi definido neste modelo.
\chapter{Introdução}
\label{cap:introducao}

A computação e a universalização do acesso a informação facilitaram e agilizaram 
em muito os processos de descoberta científica e inovação tecnológica de modo geral. 
Com isso, os problemas científicos têm se tornado cada vez mais complexos, exigindo a capacidade
de lidar com grandes volumes de dados e realizar operações computacionais de extrema complexidade que fogem 
da capacidade padrão de sistemas computacionais tradicionais. 

Para lidar com isso, infraestruturas de Computação de Alto-Desempenho 
(High-Performance Computing, HPC) têm se tornado
cada vez mais fundamentais, oferecendo altíssimos poder de processamento e capacidade de armazenamento, 
além de permitir a paralelização massiva dos algoritmos. Assim, avanços importantes vêm sendo 
feitos em campos como a biologia, onde a HPC possibilita o manejo da enorme quantidade
de dados envolvendo trabalhos com o genoma, além de possibilitar a manipulação e consequente entendimento
das estruturas e funções de proteínas em um nível molecular; como a climatologia, tornando tangível o 
desenvolvimento de modelos climáticos complexos que conseguem simular milhões de pontos atmosféricos e, através disso, 
prever possíveis desastres climáticos ou entender padrões climáticos a longo prazo com uma grande acurácia. (vale citar cybersegurança talvez)

Dada essa importância em áreas chave da ciência, infraestruturas de HPC têm cada vez mais sido tratadas
como um recurso estratégico pelas potências globais como EUA, China, Japão e Europa. Já no início dessa década, por exemplo,
os EUA inauguraram o primeiro centro de HPC público exascale, o Frontier, que possui a capacidade de efetuar mais de um quintilhão de operações
de ponto-flutuante por segundo. 

Entretanto, toda essa potência tem um preço: estima-se que o Frontier consuma cerca de 21 MW em operação contínua, 
energia que seria suficiente para abastecer cerca de 15.000 casas. Essa altíssima demanda energética é um
dos pontos chave quando se fala de HPC. Dados estimam que centros de HPC em específico consumam entre 3.2 a 5.9 bilhões 
de kWh anualmente, enquanto que datacenters de modo geral consumam cerca de 200 a 460 TWh anualmente, o que 
representa cerca de 2\% do consumo global de energia. É esperado que, num futuro próximo, essa proporção
de consumo energético se mantenha, com previsões de que o consumo deva chegar até cerca de 752 TWh, representando cerca de 2.13\% do consumo energético global. 

Esse consumo energético elevado causado pela operação dos datacenters acaba tendo um elevado
impacto ambiental, sendo notórios as emissões de gases causadores do efeito estufa (abreviados como $CO_{2}eq$) e também
o consumo de água. Diferentes métodos de geração de energia elétrica possuem diferentes pegadas
de carbono, onde, por exemplo, a fissão nuclear se caracteriza como um método com baixa emissão de $CO_{2}eq$, 
enquanto métodos baseados na queima de combustíveis fósseis como o carvão ou gás natural emitem
altas quantidades desses gases. 

Além da intensidade de carbono variar de acordo com a fonte de geração, 
temos que a composição das matrizes energéticas varia entre os países ou até mesmo entre diferentes
regiões de um mesmo país, e principalmente, varia com o passar do tempo levando em conta um dia/semana
e até mesmo de acordo com a estação do ano. Em países que tem apostado na descarbonização da rede,
podemos ter uma predominância de energia solar na matriz energética de um determinado lugar, porém 
com o chegar da noite, devido a impossibilidade da mesma, pode-se acabar queimando combustíveis
fósseis para suprimir a demanda energética.

Não obstante o impacto ambiental decorrido da operação de datacenters, temos uma fração não desprezível desses gases
incorporados no próprio processo de fabricação dos componentes dos nós dos datacenters, 
incluindo CPU, GPU, memória, dentre outros. O que agrava esse cenário é que o é possível definir
as emissões de carbono como um problema de estoque global: uma tonelada emitida em Nova York ou na China
tem o mesmo efeito quando pensamos em planeta, com o dano sendo difundido no espaço. Caso os padrões de emissão atuais 
não sejam controladas, as consequências previstas serão as mais variadas possíveis, com inundações de
cidades costeiras inteiras, grandes períodos de secas e também chuvas torrenciais, causando a disrupção
do processo agrícola, entre outros.

A outra faceta de impacto ambiental da operação dos datacenters é justamente o 
consumo de água causado por eles. Estima-se que cerca de 2/3 desse consumo decorra
indiretamente da quantidade de água consumida para gerar a energia elétrica que 
abastecem esses centros, enquanto cerca de 1/4 são consumidos nos sistemas de resfriamento
dos mesmos. Estimativas recentes atribuem cerca de 560 bilhões de litros sendo consumidos
globalmente para a operação de datacenters. 

Para piorar, boa parte dos datacenters estão situados 
em regiões com um alto indíce de estresse hídrico, acarretando em uma exploração desproporcional 
de fontes de águas em situações de estresse. Uma quantia de água consumida advinda de uma fonte
hídrica em estresse, em um período de seca ou de alto uso por exemplo, acaba competindo diretamente
com o abastecimento humano, agricultura ou outros usos societais fundamentais, tendo assim um
custo comparativamente bem maior em relação a água consumida de um fonte em estado saudável.

Apesar disso, ainda não existe uma metodologia padronizada para realizar a contabilização
da água consumida por infraestruturas de computação. A métrica de Water Usage Effectiveness (WUE),
por exemplo, leva em conta apenas a água consumida para resfriamento (on-site), deixando de 
fora a água consumida indiretamente na geração de energia utilizada. Além disso, ela se mostra incapaz
de contabilizar o fator de estresse dado fatores geográficos e sazonais da água consumida. Por fim,
boa parte dos fornecedores de Cloud não publicam métricas de eficiência hídrica, ou mesmo chegam a
publicar a quantidade de águas consumidas nas suas operações.

Com isso, boa parte dos estudos e artefatos relacionados a sustentabilidade computacional
seguem focados majoritariamente na pegada de carbono. Porém, esses dois impactos não são
sempre diretamente relacionados, a depender de fatores geográficos e temporais. A geração de energia elétrica
via usinas eólicas, por exemplo, emite uma baixa quantidade de $CO_{2}$ e consume uma baixa quantia de água,
enquanto a geração por fissão nuclear pode demandar uma grande quantia de água dependendo 
do seu método de resfriamento, enquanto emite uma baixa quantia de $CO_{2}$. Portanto, uma intervenção
poderia possívelmente reportar uma redução na sua pegada de carbono, enquanto acaba transferindo o impacto
para o consumo de água, e ainda assim registrar isso simploriamente como ganho em uma métrica agregada.
Consequentemente, avaliar qualquer intervenção em HPC com preocupação ambiental deve contabilizar ao menos
essas duas pegadas de forma isolada primeiro, e posteriormente, de forma simultânea, analisando ganhos e perdas.

\section**{O escalonador como ponto de intervenção}
\label{sec:escalonador_como_intervencao}

De modo geral, sistemas HPC de produção recebem os programas que seus usuários
desejam executar como jobs, que resumidamente, são um conjunto envolvendo o código
em si, o número de máquinas requisitadas e um tempo máximo esperado de execução. 
Esses jobs são enviados não diretamente para os sistemas operacionais dos nós dos sistemas,
mas sim para um programa intermediário, o escalonador. Portanto, é o escalonador
que decide em quais nós e quando um certo job deve rodar.

Na literatura, pode-se encontrar diversos algoritmos e heurísticas que visam otimizar certos 
aspectos na execução desses jobs, com alguns desses encontrando seu espaço nos sistemas em produção.
O EASY backfilling é um dos algoritmos mais difundidos na prática, com uma ideia geral
de simples implementação, fácil rastreio e também boa performance: o escalonador ordena a fila de
execução pelo tempo de chegada dos jobs, executando o primeiro job da fila sempre que a quantidade
de nós é suficiente. Ao mesmo tempo, o escalonador também tenta preencher possíveis nós ociosos com
jobs que estão em posições posteriores na fila, mas com uma quantidade requisitada menor ou igual
a de nós ociosos naquele momento. 

Em consequência disso, o escalonador acaba se tornando o ponto focal de intervenção quanto a 
prioridades acerca do funcionamento dos sistemas de HPC. Caso preocupações acerca do impacto ambiental
sejam o ponto focal de um certo grupo de operadores, seria possível, por exemplo, deslocar a execução
dos jobs para janelas de tempo onde a intensidade de carbono decorrente da matriz elétrica seja menor.
Essa é a premissa do escalonamento consciente de carbono, estudado majoritariamente em infraestruturas
de nuvem. Wiesner et al. reportou economias de até 20\% ao deslocar jobs tolerantes a atrasos para horários
foras de pico ou para o fim de semana, enquanto Chen et al. alegou economias de até 67\% ao combinar
deslocamentos temporais e geográficos.

A maioria dos trabalhos sobre escalonamento consciente de carbono incorporam algum tipo de previsão
como um dos seus componentes centrais nos seus algoritmos de decisão. Boa parte desses, inclusive, 
integra previsões obtidas através de modelos de Machine Learning para prever as composições das matrizes 
em um futuro próximo e poder agendar suas cargas de trabalho de acordo. Este trabalho se propõe a 
não adotar modelos preditivos treinados, optando por tratar o grau de informação sobre as intensidades
futuras como um dos seus fatores experimentais, variando-o entre a total ausência e o conhecimento perfeito,
medindo portanto quando de previsão é realmente necessária. Além disso, mesmo trabalhos com previsão
alegam ter suas reduções na pegada por meio do deslocamento temporal limitadas por fatores como baixa variabilidade local na matriz energética,
cargas de trabalho contínuas e dependentes de entradas de usuários, além da presença de cargas de trabalho
não-interrompíveis.

Uma quantia significativa de estudos sobre esse tema de escalonamento ambientalmente consciente 
argumenta que as reduções obtidas através do deslocamento geográfico dos jobs são significativamente
maiores em relação as reduções obtidas através do deslocamento temporal. Porém, quando se refere a sistemas 
HPC,nota-se que boa parte dos jobs acaba operando em um conjunto de dados significativamente maior quando
comparado com os de aplicações comumente executadas em nuvem, o que acaba por dificultar a migração
geográfica de jobs em HPC. Além disso, cada job em HPC tem uma qualidade de certa rigidez quando comparados às
aplicações em nuvem, solicitando uma quantidade exata de nós e por um certo período de tempo esperado também. Por fim,
apesar de muitos centros implementarem esquemas de preempção, o processo de salvar o estado atual (checkpoint) em jobs
MPI é custoso.

Como definido anteriormente, um job em HPC possui em sua especificação um número de nós e uma quantidade
máxima de execução necessários para rodar. Como esses são constantes durante a execução do job, e com
um modelo que considera a potência de um nó em execução como constante também, a energia ativa de um job
é demonstrada como invariante ao escalonamento. A partir disso, restam ao controle do escalonador duas
quantidades: a intensidade que multiplica a energia ativa, controlada através do deslocamento da execução do job
ao redor do tempo; e o tempo que o job espera na fila. Além disso, o adiamento na execução do job só compensa
caso seja previr uma queda na intensidade ambiental de forma que se supere o custo ocioso acumulado, ou seja,
sem informação sobre o futuro, adiar nunca compensa.

Portanto, observa-se que as duas alavancas acabam se opondo, visto que reduzir a intensidade durante 
o período ativo custa ociosidade, e a máquina ociosca durante a espera não produz nada, impossibilitada
de executar outros jobs. Com isso, é possível adicionar uma terceira variável sobre controle do escalonador:
o controle sobre a contiguidade dos jobs. Ao interromper-se um job, libera-se os seus nós para possíveis outros jobs,
convertendo energia ociosa em energia útil. Assim, adiciona-se também um terceiro preço, que é a energia de 
reexecução, termo este que é capaz de aumentar a energia ativa.

Consequentemente, conclui-se que existe um limite tangível e mensurável para o que qualquer política 
de escalonamento pode reduzir de impacto ambiental sob esses diferentes modelos. Com esse limite desconhecido,
não é possível distinguir uma política que aproveitou quase todo o espaço de otimização disponível
em relação a outra que aproveitou pouco de um espaço amplo. Por consequência, a pergunta pertinente não foca
apenas em quanto uma determinada heurística ecnomiza, mas sim em quanto há a economizar e que fração desse total
é alcançável por políticas de escalonamento que operem sem modelos preditivos treinados. 

\section**{Problema e lacuna}

Embora seja possível pesquisar por material acadêmico na Internet usando
sistemas de busca ``comuns'', existem ferramentas dedicadas, como o
\textsf{Google Scholar}\index{Google Scholar} (\url{scholar.google.com}).
O \textsf{Web of Science}\index{Web of Science} (\url{webofscience.com})
e o \textsf{Scopus}\index{Scopus} (\url{scopus.com}) oferecem recursos
sofisticados e limitam a busca a periódicos com boa reputação acadêmica.
Essas duas plataformas não são gratuitas, mas os alunos da USP têm
acesso a elas através da instituição. Algumas editoras, como a ACM
(\url{dl.acm.org}) e a IEEE (\url{ieeexplore.ieee.org}), também têm
sistemas de busca bibliográfica. Todas essas ferramentas são capazes de
exportar os dados para o formato .bib, usado pelo \LaTeX{} (no Google
Scholar, é preciso ativar a opção correspondente nas preferências).
O sítio \url{liinwww.ira.uka.de/bibliography} também permite buscar e
baixar referências bibliográficas relevantes para a área de computação.

Lamentavelmente, ainda não existe um mecanismo de verificação ou validação
das informações nessas plataformas. Portanto, é fortemente sugerido validar
todas as informações de tal forma que as entradas bib estejam corretas.
De qualquer modo, tome muito cuidado na padronização das referências
bibliográficas: ou considere TODOS os nomes dos autores por extenso, ou TODOS
os nomes dos autores abreviados. Evite misturas inapropriadas.\looseness=1

Apenas uma parte dos artigos acadêmicos de interesse está disponível livremente
na Internet; os demais são restritos a assinantes. A CAPES assina um grande
volume de publicações e disponibiliza o acesso a elas para diversas
universidades brasileiras, entre elas a USP, através do seu portal de
periódicos (\url{periodicos.capes.gov.br}). Existe uma extensão para os
navegadores Chrome e Firefox (\url{www.infis.ufu.br/capes-periodicos})
que facilita o uso cotidiano do portal.

Para manter um banco de dados organizado sobre artigos e outras fontes
bibliográficas relevantes para sua pesquisa, é altamente recomendável
que você use uma ferramenta como Zotero~(\url{zotero.org})\index{Zotero}
ou Mendeley~(\url{mendeley.com})\index{Mendeley}. Ambas podem exportar
seus dados no formato .bib, compatível com \LaTeX{}.
